\documentclass[
    language=english,
    doctype=standard,
    institution=none,
    titlestyle=book,
    oneside,
]{../../omnilatex}

\RequirePackage{config/document-settings}

\addbibresource{../../bib/bibliography.bib}

\begin{document}

\maketitle

\tableofcontents

%=============================================================================
\chapter{Introduction}
\label{sec:introduction}
%=============================================================================

This document defines the coding standards for the DataPipe project. All
contributors must follow these standards to maintain code quality,
readability, and consistency across the codebase. These standards are
enforced through automated linting, code review, and CI checks.

\section{Scope}

These standards apply to:

\begin{itemize}
    \item Java source code in the \texttt{core/}, \texttt{connectors/},
        \texttt{transforms/}, \texttt{admin-api/}, and \texttt{cli/}
        modules.
    \item Python source code in the CLI and custom transform scripts.
    \item YAML configuration files.
    \item Shell scripts in the \texttt{deploy/} directory.
\end{itemize}

%=============================================================================
\chapter{Java Code Standards}
\label{sec:java}
%=============================================================================

\section{Naming Conventions}

\begin{table}[htbp]
    \centering
    \caption{Java naming conventions.}
    \label{tab:java-naming}
    \begin{tabular}{@{}lp{5cm}p{5cm}@{}}
        \toprule
        Element & Convention & Example \\
        \midrule
        Classes & PascalCase, noun & \texttt{KafkaSourceConnector} \\
        Methods & camelCase, verb & \texttt{processRecord} \\
        Constants & SCREAMING\_SNAKE\_CASE & \texttt{MAX\_BATCH\_SIZE} \\
        Variables & camelCase & \texttt{consumerGroup} \\
        Packages & lowercase, dot-separated & \texttt{io.datapipe.core} \\
        Type parameters & Single uppercase letter & \texttt{T}, \texttt{E} \\
        \bottomrule
    \end{tabular}
\end{table}

\section{File Structure}

\begin{lstlisting}[language=java, caption={Standard Java class structure.}]
package io.datapipe.core.pipeline;

import java.util.List;
import java.util.Objects;

import io.datapipe.core.config.Config;
import io.datapipe.core.record.Record;

/**
 * Executes a pipeline by orchestrating source, transform,
 * and sink stages.
 */
public class PipelineExecutor {

    private static final int DEFAULT_BATCH_SIZE = 1000;

    private final Config config;
    private final List<Stage> stages;
    private volatile boolean running;

    public PipelineExecutor(Config config, List<Stage> stages) {
        this.config = Objects.requireNonNull(config);
        this.stages = List.copyOf(stages);
        this.running = false;
    }

    public void start() {
        if (running) {
            throw new IllegalStateException("Already running");
        }
        running = true;
        execute();
    }

    public void stop() {
        running = false;
    }

    private void execute() {
        // Implementation
    }
}
\end{lstlisting}

\section{Code Style Rules}

\begin{enumerate}
    \item Maximum line length: 100 characters.
    \item Four-space indentation (no tabs).
    \item Use \texttt{final} for all local variables that are not
        reassigned.
    \item Use \texttt{var} only when the type is obvious from context.
    \item Always use braces for \texttt{if}/\texttt{else}/\texttt{for}
        blocks, even single-line bodies.
    \item Order class members: static fields, instance fields,
        constructors, methods, inner classes.
\end{enumerate}

%=============================================================================
\chapter{Python Code Standards}
\label{sec:python}
%=============================================================================

\section{Style Guide}

Python code follows PEP 8 with the following additions:

\begin{itemize}
    \item Maximum line length: 88 characters (Black formatter default).
    \item Use type hints for all public function signatures.
    \item Use docstrings (Google style) for all public classes and
        functions.
    \item Prefer f-strings over \texttt{\%} formatting or
        \texttt{str.format}.
\end{itemize}

\section{Example}

\begin{lstlisting}[language=python, caption={Python coding standard example.}]
from __future__ import annotations

import logging
from typing import Optional

logger = logging.getLogger(__name__)


class PipelineConfig:
    """Configuration for a DataPipe pipeline.

    Attributes:
        name: Unique pipeline identifier.
        batch_size: Number of records per batch.
        timeout_ms: Maximum processing time per batch.
    """

    def __init__(
        self,
        name: str,
        batch_size: int = 1000,
        timeout_ms: int = 30000,
    ) -> None:
        self.name = name
        self.batch_size = batch_size
        self.timeout_ms = timeout_ms

    def validate(self) -> bool:
        """Validate configuration values.

        Returns:
            True if configuration is valid.

        Raises:
            ValueError: If any configuration value is invalid.
        """
        if self.batch_size <= 0:
            raise ValueError(
                f"batch_size must be positive, got {self.batch_size}"
            )
        return True
\end{lstlisting}

%=============================================================================
\chapter{File Structure}
\label{sec:file-structure}
%=============================================================================

\section{Repository Layout}

\begin{lstlisting}[language=bash, caption={Repository directory structure.}]
datapipe/
  core/
    src/main/java/io/datapipe/core/
      config/          # Configuration classes
      pipeline/        # Pipeline engine
      record/          # Record model
      transform/       # Transform interface
    src/test/java/     # Unit tests
  connectors/
    src/main/java/io/datapipe/connectors/
      kafka/           # Kafka connector
      postgresql/      # PostgreSQL connector
      s3/              # S3 connector
  transforms/
    src/main/java/io/datapipe/transforms/
      filter/          # Filter transform
      map/             # Map transform
      aggregate/       # Aggregate transform
  cli/
    src/main/python/datapipe_cli/
      commands/        # CLI commands
      utils/           # Shared utilities
  deploy/
    terraform/         # Infrastructure as code
    helm/              # Helm charts
    docker/            # Dockerfiles
  docs/                # Documentation
  scripts/             # Build and CI scripts
\end{lstlisting}

\section{File Naming}

\begin{table}[htbp]
    \centering
    \caption{File naming conventions.}
    \label{tab:file-naming}
    \begin{tabular}{@{}lp{8cm}@{}}
        \toprule
        Type & Convention \\
        \midrule
        Java class & \texttt{PascalCase.java} (matches class name) \\
        Java test & \texttt{ClassNameTest.java} \\
        Python module & \texttt{snake\_case.py} \\
        Python test & \texttt{test\_module.py} \\
        Configuration & \texttt{kebab-case.yaml} \\
        Shell script & \texttt{kebab-case.sh} \\
        \bottomrule
    \end{tabular}
\end{table}

%=============================================================================
\chapter{Code Review Checklist}
\label{sec:code-review}
%=============================================================================

All pull requests must pass the following checklist before merge:

\section{Correctness}

\begin{enumerate}
    \item Does the code do what the PR description claims?
    \item Are edge cases handled (null, empty, overflow)?
    \item Are error messages clear and actionable?
    \item Is the code thread-safe where required?
\end{enumerate}

\section{Quality}

\begin{enumerate}
    \item Is the code readable without comments explaining \emph{what} it
        does?
    \item Are variable and function names descriptive?
    \item Is there unnecessary duplication?
    \item Are magic numbers replaced with named constants?
\end{enumerate}

\section{Testing}

\begin{enumerate}
    \item Are unit tests provided for all new code?
    \item Do tests cover the happy path and error paths?
    \item Are tests independent and deterministic?
    \item Does the test name describe what it verifies?
\end{enumerate}

\section{Documentation}

\begin{enumerate}
    \item Are public APIs documented with Javadoc?
    \item Is the CHANGELOG updated for user-facing changes?
    \item Are README files updated for new features?
\end{enumerate}

%=============================================================================
\chapter{CI/CD Requirements}
\label{sec:cicd}
%=============================================================================

\section{CI Pipeline Stages}

\begin{table}[htbp]
    \centering
    \caption{CI pipeline stages and requirements.}
    \label{tab:ci-stages}
    \begin{tabular}{@{}lp{6cm}p{3cm}@{}}
        \toprule
        Stage & Description & Must Pass? \\
        \midrule
        Build & Compile all modules & Yes \\
        Unit Tests & Run unit tests with coverage & Yes \\
        Integration Tests & Run integration tests with Docker & Yes \\
        Lint & Checkstyle (Java), Ruff (Python) & Yes \\
        Security Scan & Snyk dependency scan & Yes \\
        Code Coverage & Minimum 80\% line coverage & Yes \\
        Performance & JMH benchmark regression check & No (warning) \\
        \bottomrule
    \end{tabular}
\end{table}

\section{Linting Configuration}

\begin{lstlisting}[language=bash, caption={Run linting locally.}]
# Java: Checkstyle
./gradlew checkstyleMain checkstyleTest

# Python: Ruff
ruff check cli/
ruff format --check cli/

# YAML: yamllint
yamllint deploy/**/*.yaml
\end{lstlisting}

\section{Commit Hooks}

The \texttt{pre-commit} framework runs the following checks before each
commit:

\begin{lstlisting}[language=bash, caption={pre-commit configuration.}]
# .pre-commit-config.yaml
repos:
  - repo: https://github.com/pre-commit/pre-commit-hooks
    rev: v4.5.0
    hooks:
      - id: trailing-whitespace
      - id: end-of-file-fixer
      - id: check-yaml
      - id: check-added-large-files

  - repo: https://github.com/psf/black
    rev: 24.2.0
    hooks:
      - id: black
        language_version: python3.12
\end{lstlisting}

%=============================================================================
\chapter{Enforcement}
\label{sec:enforcement}
%=============================================================================

\section{Automated Enforcement}

\begin{itemize}
    \item CI pipeline blocks merge on any failing check.
    \item Checkstyle violations are reported as PR comments.
    \item Code coverage is enforced via Codecov thresholds.
    \item Dependency vulnerabilities block merge if severity is critical.
\end{itemize}

\section{Manual Enforcement}

\begin{itemize}
    \item All PRs require at least one approval from a module owner.
    \item Large PRs (> 500 lines) require two approvals.
    \item The code review checklist must be completed.
    \item Standards exceptions require approval from the Tech Lead.
\end{itemize}

\printbibliography[heading=bibintoc, title={References}]

\end{document}
